\uuid{c5jf}
\chapitre{Statistique}
\niveau{L2}
\module{Analyse de données}
\sousChapitre{Autre}
\titre{Loi binomiale sur tableur}
\theme{statistiques, tableur}
\auteur{}
\datecreate{2022-09-30}
\organisation{AMSCC}
\difficulte{}
\contenu{
\question{Créer sur tableur la table des probabilités et des probabilités cumulées de $\mathcal B(100,0{,}4)$, puis le diagramme en bâtons. Faire dépendre les tables et le graphique des paramètres $n$ et $p$.}
\indication{Placer $n$ et $p$ dans deux cellules de paramètres, puis utiliser des références absolues.}
\reponse{Placer $n=100$ en B1 et $p=0{,}4$ en B2. En A5, utiliser \texttt{=SEQUENCE(\$B\$1+1;1;0;1)} pour obtenir automatiquement $k=0,\ldots,n$ (Excel récent).
En B5 : \texttt{=LOI.BINOMIALE.N(A5;\$B\$1;\$B\$2;FAUX)} ; en C5, la même formule avec \texttt{VRAI} donne la probabilité cumulée. Recopier sur les lignes produites.
Pour une mise à jour complète, transformer ces calculs en formules matricielles utilisant \texttt{A5\#} au lieu de A5, et définir la série du graphique avec les plages déversées \texttt{A5\#} et \texttt{B5\#}. Sans fonctions matricielles, adapter le nombre de lignes et la plage du graphique lorsque $n$ change.
Contrôles : somme des probabilités égale à 1, cumul final égal à 1, espérance 40 et variance 24. Le diagramme des probabilités est centré autour de 40 ; le cumul est croissant.}
}
